\chapter{Notes on the Sample Data}	% *NOT* \OnePageChapter

\paragraph{About appendices:}
	Each appendix follows the same page-numbering rules
	as a regular chapter; the first page of a
	(multi-page) appendix is not numbered.
	Tables and figures in an appendix are numbered
	by its letter, as Table~\ref{plottable} shows.
	The plots, measurements and notes below are
	invented for this example and describe no real study.

\begin{table}[htb]
    \caption[Invented measurements for six sample plots]{
	Invented measurements for six sample plots, included
	only to show how a table is numbered and captioned
	inside an appendix.
	}
    \begin{center}
    \begin{tabular}{||l|c|c|c|c||} \hline
	plot & area & rows & seedlings & survival \\
	     & (m$^2$) &   &           & (\%) \\  \hline \hline
	A  & 12.0 & 4 & 96  & 91 \\ \hline
	B  & 10.5 & 3 & 72  & 88 \\ \hline
	C  & 14.2 & 5 & 120 & 94 \\ \hline
	D  &  9.8 & 3 & 66  & 79 \\ \hline
	E  & 11.6 & 4 & 90  & 86 \\ \hline
	F  & 13.1 & 4 & 104 & 92 \\ \hline
	\end{tabular}
    \end{center}
\label{plottable}
\end{table}

The numbered notes that follow record how the invented
data in Table~\ref{plottable} were gathered and coded.
A list like this is a common use for an appendix: it
keeps procedural detail out of the main chapters while
leaving it available to a careful reader.

\begin{enumerate}

\item
Each plot was marked out with four corner stakes and a
length of string, and its area was measured twice. The
table reports the mean of the two measurements, rounded
to one decimal place.

\item
Rows were laid out along the longer side of each plot.
Where a plot was too narrow for a full row, the partial
row was left unplanted and is not counted.

\item
Seedlings were planted at a fixed spacing within each
row. The seedling count is the number planted on the
first day, not the number present at the end of the
season.

\item
Survival is the share of planted seedlings still alive
at the final visit, expressed as a whole-number
percentage.

\item
Visits took place once a week. A visit that was missed
because of weather was made up on the next dry day and
recorded under its original week.

\item
A seedling was counted as alive if it had at least one
green leaf. Seedlings that had been uprooted but were
replanted on the same visit were counted as alive.

\item
Plot~D lies closest to the path and lost several
seedlings to foot traffic early in the season. No
adjustment was made for this in the table.

\item
Watering followed the same schedule in every plot. On
the two hottest weeks an extra watering was added to all
six plots at once, so the schedule stayed the same
across plots.

\item
No fertilizer was applied. A single layer of mulch was
spread over every plot after the third visit.

\item
Weeds were removed by hand at each visit. The time spent
weeding was noted but is not reported here.

\item
Measurements were written in a field notebook and typed
into a spreadsheet on the same evening. Each typed entry
was checked against the notebook by a second reader.

\item
Two entries were corrected during that check: one
transposed digit in the area of Plot~B and one miscounted
row in Plot~E. The table shows the corrected values.

\item
All numbers in this appendix are fictional. They were
chosen to fill a table of realistic size, and no
conclusion should be drawn from them.

\item
Readers adapting this example for their own thesis
should replace this appendix entirely with their own
material, keeping only the structure if it is useful.

\end{enumerate}
